% erzeugt von bau.py aus zone_a, e1_a, e2_a, e3_a
\begin{abhakseite}
\abhakgruppe{Kennst du schon}
\abhak{1}{Ich kann mit Kommazahlen malnehmen und teilen und das Ergebnis runden.}
\abhak{2}{Ich kann Umfang und Flächeninhalt unterscheiden und mit der passenden Einheit angeben.}
\abhak{3}{Ich kann eine Formel nach einer anderen Größe umstellen.}
\abhak{4}{Ich kann Zahlen quadrieren und die Wurzel ziehen.}
\abhak{5}{Ich finde den Fehler beim Quadrieren und Wurzelziehen.}
\abhak{6}{Ich kann Quadrat und Doppeltes auseinanderhalten.}
\abhak{7}{Ich kann Winkel messen und zum Vollwinkel ergänzen.}
\abhak{8}{Ich kann einen Winkel mit dem Geodreieck zeichnen.}
\abhak{9}{Ich kann einen Anteil als Bruch und in Prozent angeben.}
\abhakgruppe{Einheit 1 · Kreisumfang}
\abhak{10}{Ich erkenne, ob Radius oder Durchmesser gegeben ist.}
\abhak{11}{Ich erkenne, ob der Rand oder die Fläche gefragt ist.}
\abhakgruppe{\hspace*{1em}\mdseries\itshape Radius und Durchmesser}
\abhak{12}{Ich kann den Durchmesser aus dem Radius berechnen und umgekehrt.}
\abhakgruppe{\hspace*{1em}\mdseries\itshape Kreise zeichnen}
\abhak{13}{Ich kann Mittelpunkt, Radius und Durchmesser in einen Kreis einzeichnen (kennst du seit Klasse 5).}
\abhak{14}{Ich kann einen Kreis mit dem Zirkel zeichnen.}
\abhakgruppe{\hspace*{1em}\mdseries\itshape Umfang und Durchmesser vergleichen}
\abhak{15}{Ich kann Umfang und Durchmesser eines Kreises vergleichen.}
\abhakgruppe{\hspace*{1em}\mdseries\itshape Umfang berechnen}
\abhak{16}{Ich kann den Umfang eines Kreises berechnen.}
\abhak{17}{Ich kann aus dem Umfang den Durchmesser und den Radius berechnen.}
\abhak{18}{Ich kann die Länge von Kreisrändern in Figuren und Körpern berechnen.}
\abhak{19}{Ich finde den Fehler in einer Umfangsrechnung.}
\abhak{20}{Ich kann begründen, warum Umfang geteilt durch Durchmesser bei jedem Kreis gleich ist.}
\abhak{21}{Ich kann ausrechnen, wie weit ein Rad rollt.}
\abhakgruppe{Einheit 2 · Kreisfläche}
\abhak{22}{Ich erkenne, welche Formel passt.}
\abhakgruppe{\hspace*{1em}\mdseries\itshape Fläche aus Radius oder Durchmesser}
\abhak{23}{Ich kann die Fläche eines Kreises berechnen.}
\abhak{24}{Ich kann die Fläche von Halbkreis und Viertelkreis berechnen.}
\abhakgruppe{\hspace*{1em}\mdseries\itshape Flächenterme zu Kreisfiguren}
\abhak{25}{Ich kann zu einer Kreisfigur den Term für die Fläche aufschreiben.}
\abhak{26}{Ich kann zu einem Flächenterm die Kreisfigur angeben.}
\abhakgruppe{\hspace*{1em}\mdseries\itshape Radius, Durchmesser, Umfang und Fläche}
\abhak{27}{Ich kann in einer Tabelle Radius, Durchmesser, Umfang und Fläche ergänzen.}
\abhak{28}{Ich kann zwischen Fläche, Radius und Durchmesser umrechnen.}
\abhak{29}{Ich finde den Fehler in einer Flächenrechnung.}
\abhak{30}{Ich kann begründen, wie die Fläche vom Radius abhängt.}
\abhak{31}{Ich kann unterscheiden, ob Umfang oder Fläche gesucht ist.}
\abhak{32}{Ich kann mit Kreisflächen Preise vergleichen.}
\abhakgruppe{Einheit 3 · Kreisteile}
\abhak{33}{Ich erkenne, welcher Teil vom Kreis gemeint ist.}
\abhakgruppe{\hspace*{1em}\mdseries\itshape Anteil und Mittelpunktswinkel}
\abhak{34}{Ich kann den Anteil eines Kreisausschnitts am ganzen Kreis angeben.}
\abhak{35}{Ich kann aus dem Anteil den Mittelpunktswinkel berechnen.}
\abhakgruppe{\hspace*{1em}\mdseries\itshape Bogen, Ausschnitt und Ring}
\abhak{36}{Ich kann die Länge eines Kreisbogens berechnen.}
\abhak{37}{Ich kann die Fläche eines Kreisausschnitts berechnen.}
\abhak{38}{Ich kann den Umfang eines Kreisausschnitts berechnen.}
\abhak{39}{Ich kann die Fläche eines Kreisrings berechnen.}
\abhak{40}{Ich finde den Fehler bei Kreisteilen.}
\abhak{41}{Ich kann begründen, wie Mittelpunktswinkel und Anteil zusammenhängen.}
\abhak{42}{Ich kann mit einem Kreisausschnitt eine Rasenfläche berechnen.}
\abhak{43}{Ich kann ein Kreisdiagramm berechnen und zeichnen.}
\end{abhakseite}
